\documentclass[10pt,a4paper]{amsart}
\usepackage[margin=19mm]{geometry}
\usepackage{amsmath,amssymb,mathtools}
\usepackage[scr=rsfs,cal=euler]{mathalfa}
\usepackage{xfrac}
\usepackage[unicode,hidelinks,bookmarks=false]{hyperref}
\newcommand{\ie}{i.e.\ }
\newcommand{\cf}{cf.\ }
\newcommand{\resp}{respectively\ }
\newcommand{\Subsection}{\subsection}
\providecommand{\nobreakdash}{\nobreak\hbox{-}}
\setlength{\emergencystretch}{2em}
\multlinegap0pt
\usepackage{bm}
\usepackage{bbm}
\usepackage{dsfont}
\usepackage{xspace}
\usepackage{cancel}
\usepackage{mathtools}
\usepackage{amsthm}
\makeatletter
\@ifundefined{theorem}{\newtheorem{theorem}{Theorem}}{}
\@ifundefined{lemma}{\newtheorem{lemma}[theorem]{Lemma}}{}
\makeatother
\title[On forbidden configurations in point-line incidence graphs]{On forbidden configurations in point-line incidence graphs}
\author{Martin Balko, N\'ora Frankl}
\date{Source: arXiv:2409.00954v2 (CC BY 4.0)}
\begin{document}
\maketitle

\noindent\emph{Source and licence.} On forbidden configurations in point-line incidence graphs. Version arXiv:2409.00954v2 (CC BY 4.0), distributed under the Creative Commons Attribution 4.0 International licence, as recorded in the arXiv licence metadata for this exact version and shown on the abstract page https://arxiv.org/abs/2409.00954v2. Full rights notice: licenses/arxiv-2409.00954-CC-BY-4.0.txt.\par\smallskip

\noindent\textit{Editorial scope.} Counted unit: the Lemma whose source label is \texttt{lem-embedding} in the article (source0.tex lines 292--300, source label \texttt{lem-embedding}), delivered together with the article's own complete original proof of it (source0.tex lines 317--365, one \texttt{proof} environment of 49 lines). No mathematics is written, repaired or re-derived: statement and proof are the article's own text line for line. The proof is detached from the statement in the source: the article states the result at lines 292--300 and prints its proof at lines 317--365, and the proof environment's own header names the statement, so the association is the article's own. Required context retained uncounted: none. Every definition, display and label of those lines is retained unchanged. Omitted from this package and not redistributed: 1--291, 301--316, 366--691. Nothing inside a retained excerpt is omitted or altered: every retained body is delivered exactly as the article's lines are written, so no omission receipt of any kind accompanies this package. Citations: the retained excerpts carry no citation command at all, so no bibliography entry of the article is transcribed and no reference is redistributed. Reference adaptation: none was needed and none was made; every reference inside the delivered unit resolves inside it, so no unresolved reference or citation is produced. Printed slips kept literal: no printed word, symbol, hypothesis or number of the counted statement or of its proof was altered. Layout and class: the article's own class is \texttt{article} with packages including graphicx, amsthm, amssymb, amsmath, xcolor, enumerate, microtype, natbib, hyperref; the wrapper is the lane's pinned \texttt{amsart} editorial wrapper at 10pt on A4, which restates the theorem-like environments the retained text needs and adds the article's own macro definitions that the retained text needs (none), with extra wrapper packages (bm, bbm, dsfont, xspace, cancel, mathtools). The delivered numbering is the unit's own. Assets: no retained span contains a figure environment, a graphics-inclusion command, a PDF-page inclusion, a table or a TikZ picture; no image file is copied into this package, the package declares no asset dependency and no image file is redistributed with it. Licence: version arXiv:2409.00954v2 of this article is distributed under the Creative Commons Attribution 4.0 International licence, as recorded in the arXiv licence metadata for this exact version and shown on the abstract page https://arxiv.org/abs/2409.00954v2. A full rights notice is delivered at licenses/arxiv-2409.00954-CC-BY-4.0.txt. Characters: every retained excerpt is pure ASCII, so the delivered engine, XeLaTeX with its default Latin Modern text font, reports no missing character.
\begin{lemma}
\label{lem-embedding}
There is an absolute constant $c \geq 1$ such that, for all integers $d \geq 1$ and $k \geq 5$ satisfying
\begin{equation}
\label{eq-condition}
deg\left(\left(1+2\cos\left(\frac{2\pi}{k}\right)\right)^2 \right) > d^c
\end{equation}
and for every projective transformation $f$, if $R'$ is an extended regular $k$-gon, then $f(R')$ is not a subset of $V_d$.
\end{lemma}

\begin{proof}[Proof of Lemma~\ref{lem-embedding}]
For $d \geq 1$, let $R'$ be an extended regular $k$-gon with $k \geq 5$, and let $f$ be a projective transformation.
We may assume that $f(R') \subseteq \mathbb{R}^2$ as otherwise some points of $f(R')$ are mapped to the line in infinity and are not contained in~$V_d$.
As before, we use $R$ to denote the regular $k$-gon on vertices $v_1,\dots,v_k$ that is a subset of $R'$.
We let $w$ be the intersection point of the lines $\overline{v_1v_2}$ and $\overline{v_{k-1},v_k}$.
Since $k \geq 5$, the point $w$ is not one of the points from $\{v_1,\dots,v_k\}$ nor does it lie in infinity.

Let $s$ be the length of each side of $R$, that is, $s=|v_iv_{i+1}|$ for every $i \in [k]$.
We also use $r$ to denote the distance $|wv_1|$.
The points $w$, $v_1$, and $v_k$ span an isosceles triangle with sides of lengths $r$, $r$, and $s$.
Now, the angles $wv_1v_k$ and $wv_kv_1$ are equal to $(2\pi - 2\frac{(k-2)\pi}{k})/2=2\pi/k$, since the internal angle of $R$ is $\frac{(k-2)\pi}{k}$.
It follows that the angle $v_1wv_k$ equals $\pi - 2 \cdot 2\pi/k = \frac{(k-4)\pi}{k}$.
By the \emph{Law of sines}, we get
\begin{equation}
\label{eq-ratio}
\frac{s}{r} = \frac{\sin\left(\frac{(k-4)\pi}{k}\right)}{\sin\left(\frac{2\pi}{k}\right)} = \frac{\sin\left(\frac{4\pi}{k}\right)}{\sin\left(\frac{2\pi}{k}\right)} = 2\cos\left(\frac{2\pi}{k}\right),
\end{equation}
where we also used the \emph{Double angle identity} $\sin(2x) = 2\sin(x)\cos(x)$.

To compute the cross-ratio, note that, by~\eqref{eq-ratio},
\begin{equation}
\label{eq-crossRatio}
(t_1,w;v_1,v_2) = \frac{|t_1v_1||wv_2|}{|t_1v_2||wv_1|} = \frac{|wv_2|}{|wv_1|} = \frac{r+s}{r} = 1+2\cos\left(\frac{2\pi}{k}\right),
\end{equation}
where we used the fact that $t_1$ is the line's point at infinity.

Now, suppose for contradiction that $R'$ is an extended regular $k$-gon satisfying $f(R') \subseteq V_d$ and $deg\left(\left(1+2\cos\left(\frac{2\pi}{k}\right)\right)^2\right) > d^c$.

We show that the point $w$ has coordinates of degree at most $d^{c'}$ for some absolute constant $c' \geq 1$.
The point $w$ is the intersection point of two lines $\overline{v_1v_2}$ and $\overline{v_{k-1},v_k}$. 
Thus, the coordinates of $w$ can be expressed in terms of the coordinates of the points $v_1,v_2,v_{k-1},v_k$ using a fixed number of additions, multiplications, and divisions.
If $\alpha$ and $\beta$ are two algebraic numbers of degrees $m$ and $n$, respectively, then the degree of $\alpha + \beta$, $\alpha \cdot \beta$, and $\alpha^{-1}$ are at most $mn$, $mn$, and $m$, respectively.
Since $v_1,v_2,v_{k-1},v_k \in V_d$, the coordinates of all points $v_1,v_2,v_{k-1},v_k$ have degree at most $d$.
It follows that the degrees of the coordinates of $w$ are at most $d^{c'}$ for some absolute constant $c' \geq 1$.

Using the bounds on degrees of numbers obtained by additions and multiplications and the facts $f(R') \subseteq V_d$ and $w \in V_{d^c}$, the squares of all distances between two points from $f(R' \cup \{w\})$ are numbers of degree at most $d^c$ for some absolute constant $c \geq 1$.
Thus, the number $(f(t_1),f(w);f(v_1),f(v_2))^2$ has degree at most $d^c$.
Since projective transformations preserve cross-ratios, we have
\begin{equation}
\label{eq-preserving}
(t_1,w;v_1,v_2) = (f(t_1),f(w);f(v_1),f(v_2)).
\end{equation}
On the other hand, the equation~\eqref{eq-crossRatio} 
gives
\[
deg\left((t_1,w;v_1,v_2)^2\right) = deg\left(\left(1+2\cos\left(\frac{2\pi}{k}\right)\right)^2\right) > d^c.
\]
By combining this with~\eqref{eq-preserving}, we obtain that $(f(t_1),f(w);f(v_1),f(v_2))^2$ has degree larger than $d^c$, a contradiction.
\end{proof}

\end{document}
